\documentclass[11pt,a4paper]{article}
% Shared preamble for the qi26_25 weekly deliverable reports.
%
% Each weekly report is a short standalone document that inputs this file, so
% that formatting stays consistent across all six without duplicating the
% package list in every one. Change something here and every report follows.
%
% Usage in a weekly report:
%     \documentclass[11pt,a4paper]{article}
%     % Shared preamble for the qi26_25 weekly deliverable reports.
%
% Each weekly report is a short standalone document that inputs this file, so
% that formatting stays consistent across all six without duplicating the
% package list in every one. Change something here and every report follows.
%
% Usage in a weekly report:
%     \documentclass[11pt,a4paper]{article}
%     % Shared preamble for the qi26_25 weekly deliverable reports.
%
% Each weekly report is a short standalone document that inputs this file, so
% that formatting stays consistent across all six without duplicating the
% package list in every one. Change something here and every report follows.
%
% Usage in a weekly report:
%     \documentclass[11pt,a4paper]{article}
%     \input{weekly_preamble}
%     \weekhead{1}{Vision Foundation and Data Encoding}{Week of ...}
%     ... body ...
%     \end{document}

\usepackage[utf8]{inputenc}
\usepackage[T1]{fontenc}
\usepackage[margin=2.5cm]{geometry}
\usepackage{amsmath,amssymb}
\usepackage{booktabs}
\usepackage{graphicx}
\usepackage{enumitem}
\usepackage{listings}
\usepackage{xcolor}
\usepackage[hidelinks]{hyperref}
\usepackage{titlesec}
\usepackage{fancyhdr}

\definecolor{codebg}{RGB}{247,247,250}
\definecolor{codekw}{RGB}{20,90,160}
\definecolor{codecm}{RGB}{110,120,130}
\definecolor{accent}{RGB}{20,90,160}

\lstset{
  backgroundcolor=\color{codebg},
  basicstyle=\ttfamily\footnotesize,
  keywordstyle=\color{codekw}\bfseries,
  commentstyle=\color{codecm}\itshape,
  breaklines=true,
  frame=single,
  rulecolor=\color{gray!30},
  language=Python,
  showstringspaces=false,
  xleftmargin=6pt,
  framexleftmargin=6pt
}

\setlist{itemsep=2pt,topsep=4pt}
\titleformat{\section}{\normalfont\large\bfseries\color{accent}}{\thesection}{0.8em}{}
\titleformat{\subsection}{\normalfont\normalsize\bfseries}{\thesubsection}{0.8em}{}

\pagestyle{fancy}
\fancyhf{}
\fancyfoot[C]{\small\thepage}
\fancyfoot[R]{\scriptsize\texttt{github.com/Pushkar0997/qi26\_25}}
\renewcommand{\headrulewidth}{0pt}

% \weekhead{number}{title}{deliverable as stated in the project brief}
\newcommand{\weekhead}[3]{%
  \begin{center}
    {\scriptsize\textsf{QINTERN 2026 \textbullet\ QWORLD \textbullet\ WEEKLY DELIVERABLE}}\\[0.25cm]
    {\LARGE\bfseries Week #1: #2}\\[0.35cm]
    {\normalsize Hybrid QML and Quantum String Algorithms for Document Extraction}\\[0.25cm]
    {\small Pushkar Kumar \quad\textbullet\quad Mentor: Potluri Krishna Priyatham}\\[0.4cm]
    \fbox{\begin{minipage}{0.88\textwidth}
      \small\textbf{Deliverable as specified in the project brief:} #3
    \end{minipage}}
  \end{center}
  \vspace{0.4cm}
}

% A boxed note for caveats and honest limitations.
\newcommand{\honestnote}[1]{%
  \vspace{0.2cm}
  \noindent\fcolorbox{gray!40}{gray!8}{%
    \begin{minipage}{0.97\textwidth}\small #1\end{minipage}}
  \vspace{0.2cm}
}

%     \weekhead{1}{Vision Foundation and Data Encoding}{Week of ...}
%     ... body ...
%     \end{document}

\usepackage[utf8]{inputenc}
\usepackage[T1]{fontenc}
\usepackage[margin=2.5cm]{geometry}
\usepackage{amsmath,amssymb}
\usepackage{booktabs}
\usepackage{graphicx}
\usepackage{enumitem}
\usepackage{listings}
\usepackage{xcolor}
\usepackage[hidelinks]{hyperref}
\usepackage{titlesec}
\usepackage{fancyhdr}

\definecolor{codebg}{RGB}{247,247,250}
\definecolor{codekw}{RGB}{20,90,160}
\definecolor{codecm}{RGB}{110,120,130}
\definecolor{accent}{RGB}{20,90,160}

\lstset{
  backgroundcolor=\color{codebg},
  basicstyle=\ttfamily\footnotesize,
  keywordstyle=\color{codekw}\bfseries,
  commentstyle=\color{codecm}\itshape,
  breaklines=true,
  frame=single,
  rulecolor=\color{gray!30},
  language=Python,
  showstringspaces=false,
  xleftmargin=6pt,
  framexleftmargin=6pt
}

\setlist{itemsep=2pt,topsep=4pt}
\titleformat{\section}{\normalfont\large\bfseries\color{accent}}{\thesection}{0.8em}{}
\titleformat{\subsection}{\normalfont\normalsize\bfseries}{\thesubsection}{0.8em}{}

\pagestyle{fancy}
\fancyhf{}
\fancyfoot[C]{\small\thepage}
\fancyfoot[R]{\scriptsize\texttt{github.com/Pushkar0997/qi26\_25}}
\renewcommand{\headrulewidth}{0pt}

% \weekhead{number}{title}{deliverable as stated in the project brief}
\newcommand{\weekhead}[3]{%
  \begin{center}
    {\scriptsize\textsf{QINTERN 2026 \textbullet\ QWORLD \textbullet\ WEEKLY DELIVERABLE}}\\[0.25cm]
    {\LARGE\bfseries Week #1: #2}\\[0.35cm]
    {\normalsize Hybrid QML and Quantum String Algorithms for Document Extraction}\\[0.25cm]
    {\small Pushkar Kumar \quad\textbullet\quad Mentor: Potluri Krishna Priyatham}\\[0.4cm]
    \fbox{\begin{minipage}{0.88\textwidth}
      \small\textbf{Deliverable as specified in the project brief:} #3
    \end{minipage}}
  \end{center}
  \vspace{0.4cm}
}

% A boxed note for caveats and honest limitations.
\newcommand{\honestnote}[1]{%
  \vspace{0.2cm}
  \noindent\fcolorbox{gray!40}{gray!8}{%
    \begin{minipage}{0.97\textwidth}\small #1\end{minipage}}
  \vspace{0.2cm}
}

%     \weekhead{1}{Vision Foundation and Data Encoding}{Week of ...}
%     ... body ...
%     \end{document}

\usepackage[utf8]{inputenc}
\usepackage[T1]{fontenc}
\usepackage[margin=2.5cm]{geometry}
\usepackage{amsmath,amssymb}
\usepackage{booktabs}
\usepackage{graphicx}
\usepackage{enumitem}
\usepackage{listings}
\usepackage{xcolor}
\usepackage[hidelinks]{hyperref}
\usepackage{titlesec}
\usepackage{fancyhdr}

\definecolor{codebg}{RGB}{247,247,250}
\definecolor{codekw}{RGB}{20,90,160}
\definecolor{codecm}{RGB}{110,120,130}
\definecolor{accent}{RGB}{20,90,160}

\lstset{
  backgroundcolor=\color{codebg},
  basicstyle=\ttfamily\footnotesize,
  keywordstyle=\color{codekw}\bfseries,
  commentstyle=\color{codecm}\itshape,
  breaklines=true,
  frame=single,
  rulecolor=\color{gray!30},
  language=Python,
  showstringspaces=false,
  xleftmargin=6pt,
  framexleftmargin=6pt
}

\setlist{itemsep=2pt,topsep=4pt}
\titleformat{\section}{\normalfont\large\bfseries\color{accent}}{\thesection}{0.8em}{}
\titleformat{\subsection}{\normalfont\normalsize\bfseries}{\thesubsection}{0.8em}{}

\pagestyle{fancy}
\fancyhf{}
\fancyfoot[C]{\small\thepage}
\fancyfoot[R]{\scriptsize\texttt{github.com/Pushkar0997/qi26\_25}}
\renewcommand{\headrulewidth}{0pt}

% \weekhead{number}{title}{deliverable as stated in the project brief}
\newcommand{\weekhead}[3]{%
  \begin{center}
    {\scriptsize\textsf{QINTERN 2026 \textbullet\ QWORLD \textbullet\ WEEKLY DELIVERABLE}}\\[0.25cm]
    {\LARGE\bfseries Week #1: #2}\\[0.35cm]
    {\normalsize Hybrid QML and Quantum String Algorithms for Document Extraction}\\[0.25cm]
    {\small Pushkar Kumar \quad\textbullet\quad Mentor: Potluri Krishna Priyatham}\\[0.4cm]
    \fbox{\begin{minipage}{0.88\textwidth}
      \small\textbf{Deliverable as specified in the project brief:} #3
    \end{minipage}}
  \end{center}
  \vspace{0.4cm}
}

% A boxed note for caveats and honest limitations.
\newcommand{\honestnote}[1]{%
  \vspace{0.2cm}
  \noindent\fcolorbox{gray!40}{gray!8}{%
    \begin{minipage}{0.97\textwidth}\small #1\end{minipage}}
  \vspace{0.2cm}
}


\begin{document}
\weekhead{5}{Algorithmic Refinement and Data Cleaning}
{Cleaned data output ready for industrial application.}

\section{Status: Not Delivered}

This deliverable was not completed. This report states what was specified, why it
was cut, what would have been required, and what was done instead with the time.

A documented account of an attempted and abandoned objective is a legitimate part
of a research report; silently omitting the week would not be.

\section{What Was Specified}

The project brief set out three items for this week:

\begin{itemize}
\item Implement quantum subroutines for string alignment and cleaning of OCR
      noise.
\item Replace LLM-based parsing with targeted quantum logic for specific data
      schemas.
\item Deliver cleaned data output ready for industrial application.
\end{itemize}

\section{Why It Was Cut}

\subsection{The Precondition Was Not Met}

Quantum string alignment operates on extracted text. Its purpose is to correct
residual OCR errors---aligning a noisy read against a schema or a reference
vocabulary and repairing mismatches.

That presupposes an extraction stage whose output is mostly correct, with sparse
and localised errors. The measured reality at the point this week began:

\begin{table}[h]
\centering
\caption{Extraction quality available as input to a cleaning stage.}
\begin{tabular}{@{}lrr@{}}
\toprule
Tier & CER & Identifier recovered exactly \\
\midrule
clean\_digital & 6.6\% & 100\% \\
clean\_scan & 8.5\% & 58\% \\
noisy\_scan & 79.5\% & 0\% \\
clear\_handwriting & 51.8\% & 0\% \\
degraded\_handwriting & 85.7\% & 0\% \\
\bottomrule
\end{tabular}
\end{table}

On the two clean tiers, cleaning has little to work with: 100\% of clean digital
identifiers are already exact, so a cleaning stage could improve nothing there.
On the three degraded tiers, error rates of 52--86\% mean the input is not noisy
text but largely wrong text. Alignment against a schema does not repair a string
where most characters are incorrect; it produces a confident and wrong result.

\subsection{The Resource Cost Was Prohibitive}

Quantum sequence alignment via a Grover-based approach would extend the
comparator oracle of Week 4. That oracle's measured cost is the obstacle:

\begin{itemize}
\item 3{,}272 CX gates for a 16-character field with a 2-character pattern.
\item Measured scaling exponent of $1.39$ in text length---superlinear, because
      each candidate position requires its own controlled data load.
\item A full 38-symbol alphabet needs 6 qubits per character, placing even short
      alignment windows beyond feasible simulation.
\end{itemize}

Alignment requires comparing against multiple reference candidates rather than a
single pattern, multiplying an already superlinear cost. Week 4 had already
established that the $\sqrt{N}$ advantage does not survive data loading for
single-pattern search; there was no basis to expect it to survive for alignment.

\subsection{Time}

The project ran on a compressed schedule. Section~\ref{sec:scope} records the
circumstances. Given a finite budget, the choice was between attempting this
deliverable and completing the benchmarking that Week 6 specified.

\honestnote{%
Attempting this week would have produced a component that did not work, resting
on an extraction stage that could not support it, at a resource cost already
measured as prohibitive. A broken feature discovered at submission is worse than
a documented decision not to build it.}

\section{What Would Have Been Required}

For completeness, the path to delivering this week properly:

\begin{enumerate}
\item \textbf{Fix the extraction stage first.} The measured bottleneck is the
      classical segmentation front end, not the quantum stages. Clean-tier CER of
      6.6--8.5\% against 93--95\% segmentation recovery means the residual error
      is missed and merged glyphs. A learned detector replacing projection
      profiles is the correct first step, and would raise extraction quality to
      the point where cleaning has something to act on.
\item \textbf{Adopt a coarser encoding for search.} The
      one-qubit-per-candidate-position encoding is what makes the oracle
      expensive. Stem-level or block-level encodings, as used in related QUBO
      formulations, reduce both qubit count and loading cost.
\item \textbf{Define the schema constraint explicitly.} ``Targeted quantum logic
      for specific data schemas'' requires a specified schema. The synthetic
      corpus uses a fixed identifier format, which would be a reasonable starting
      constraint, but the brief did not specify a target schema and none was
      agreed.
\item \textbf{Establish a classical baseline first.} Fuzzy matching against a
      schema is a solved classical problem. Without measuring that baseline,
      there would be no way to tell whether a quantum implementation contributed
      anything---which is exactly the failure mode Weeks 2 and 4 were designed to
      avoid.
\end{enumerate}

\section{What Was Done Instead}

The time was spent on work not specified in the brief but which strengthened the
project's core claims:

\begin{itemize}
\item \textbf{Statistical significance testing.} The Week 2 comparison rested on
      a single train--test split. Repeating over 30 paired splits established
      that the quantum-versus-classical result is a null result
      ($p \approx 0.2$--$0.3$) while the raw-pixel lead is real
      ($p < 10^{-10}$). Without this the two claims were indistinguishable.
\item \textbf{An external baseline.} Tesseract was run on identical images,
      providing the reference point that the internal comparisons lacked.
\item \textbf{Noise and hardware-feasibility analysis.} Gate error, readout error
      and shot noise were modelled, establishing that the feature extractor is
      plausible on near-term hardware while the search stage is not.
\item \textbf{Reproducibility infrastructure.} Committed dataset and model,
      browser-parity verification, and continuous integration that fails when the
      report and measured results disagree.
\end{itemize}

\section{Scope}
\label{sec:scope}

The project was scoped for three participants, with separate owners for the
vision and search tracks. Mentor contact ended in early July, and the two track
assignments were not taken up. Both tracks were consolidated into a single-person
scope.

This is recorded as context for the schedule, not as justification for the
decision above. The precondition and resource arguments of Section 3 stand
independently: this deliverable would have been the correct one to cut even with
a full team.

\section{Deliverable Status}

\textbf{Not delivered.} Cut deliberately, with reasons recorded here and in
Section~8 of the final report.

\end{document}
